\section*{Introducción}

Esta conversación parece, en la superficie, tratar sobre Zhao Chenyang (赵晨阳): licenciatura en Tsinghua, quince meses de doctorado en UCLA, participación en SGLang, y el paso de un proyecto de código abierto a fundar una empresa de motores de inferencia; cuando se publicó el video, RadixArk acababa de anunciar una ronda semilla de cien millones de dólares. Si uno se queda solo en los titulares, es fácil entenderla como «la historia de un joven más en la ola de abandonos universitarios de Silicon Valley». Pero lo que realmente vale la pena conservar de toda la charla no es el abandono en sí, sino la pregunta más profunda a la que apuntan juntos el abandono, SGLang, AI Coding y la comunidad de código abierto: una vez que la IA reduce con rapidez el costo de ejecución, ¿qué responsabilidades le quedan a la persona que no puede delegar?

Zhao Chenyang insiste una y otra vez en que «hay que hacer las cosas de manera esencial». Esta idea recorre toda la charla. Cuando habla del mundo académico, critica la inercia de tener un martillo y salir a buscar clavos, de formular una evaluation en el vacío; cuando habla de dejar el doctorado, le preocupa si la applied research realmente necesita someterse a la validación del mercado y de la ingeniería; cuando habla de AI Coding, subraya que, una vez que producir código se vuelve barato, las personas necesitan hacerse cargo del design, del review, de decir que no (say no) y de asumir las consecuencias; cuando habla de la comunidad de código abierto, lo que en realidad le importa no es «cuántas contribuciones» se hacen, sino cómo un proyecto mantiene sus límites de responsabilidad, su sentido estético y su orden de colaboración una vez que el código generado por IA se vuelve abundante.

Por eso, estas notas no siguen el orden cronológico, sino que se reorganizan en torno a seis temas: la posición ecológica de los motores de inferencia, el cambio de paradigma de investigación detrás de la ola de abandonos, el research taste en la era de la IA, el impacto de AI Coding sobre la gobernanza del código, el riesgo cognitivo que trae la IA como amplificador, y cómo la comunidad de código abierto redefine la contribución en la era del código generado.

\begin{importantbox}{Tres hilos conductores para leer estas notas}
El primer hilo es la «esencialidad»: la applied research necesita una application, y la fundamental research necesita un sistema de problemas coherente consigo mismo; no basta con optimizar contra un blanco ficticio. El segundo hilo es el «límite de responsabilidad»: la IA abarata la ejecución, pero también facilita que se multipliquen los pull requests, el código y el contenido hechos sin pensar. El tercer hilo es la «comunidad y la organización»: un proyecto de código abierto no es solo un repositorio de código, sino que debe registrar contribuciones, filtrar responsabilidades, transmitir información y organizar a las personas para que trabajen juntas.
\end{importantbox}

El video original no tiene subtítulos oficiales; esta transcripción se elaboró a partir de una transcripción local con Whisper y una revisión párrafo por párrafo. El cuerpo del texto no reproduce la transcripción palabra por palabra, sino que busca extraer, en la medida de lo posible, las cadenas de juicio y los puntos de vista transferibles de la entrevista.
